\documentclass[11pt, a4paper]{article}

\usepackage{lmodern}
\usepackage{microtype}
\usepackage[margin=2.5cm]{geometry}
\usepackage[dvipsnames, table]{xcolor}
\definecolor{linkblue}{rgb}{0.0, 0.2, 0.6}
\definecolor{linkred}{rgb}{0.7, 0.13, 0.13}
\usepackage[style=numeric, backend=biber]{biblatex}
\usepackage{citemsp}
\usepackage[colorlinks=true, citecolor=linkblue,
            linkcolor=linkred, urlcolor=linkblue]{hyperref}
\usepackage{tcolorbox}
\usepackage{authblk}
\usepackage{graphicx}
\usepackage{array}
\usepackage{booktabs}
\usepackage{orcidlink}
\usepackage{fontawesome5}
\usepackage{titlesec}

% ── Section formatting ────────────────────────────────────────────────────────
\titleformat{\section}
  {\Large\sffamily\bfseries\color{black!80}}
  {\colorbox{black!10}{\makebox[1.2cm]{\thesection}}}
  {1em}
  {}
  [\vspace{-0.3em}{\color{black!20}\rule{\textwidth}{0.5pt}}]

\titleformat{\subsection}
  {\large\sffamily\bfseries\color{black!70}}
  {\thesubsection}
  {1em}
  {}

% ── Helpers ───────────────────────────────────────────────────────────────────
\newcommand{\githublink}[1]{\href{https://github.com/#1}{\faGithub}}

% ── Example bibliography ─────────────────────────────────────────────────────
\begin{filecontents}{refs-citemsp.bib}
@article{einstein1915,
  author  = {Einstein, Albert},
  title   = {Die Feldgleichungen der Gravitation},
  journal = {Sitzungsberichte der Preussischen Akademie der Wissenschaften},
  year    = {1915},
  pages   = {844--847}
}
@book{wald1984,
  author    = {Wald, Robert M.},
  title     = {General Relativity},
  publisher = {University of Chicago Press},
  year      = {1984}
}
@article{hawking1974,
  author  = {Hawking, Stephen W.},
  title   = {Black hole explosions?},
  journal = {Nature},
  volume  = {248},
  pages   = {30--31},
  year    = {1974}
}
@book{mtw1973,
  author    = {Misner, Charles W. and Thorne, Kip S. and Wheeler, John Archibald},
  title     = {Gravitation},
  publisher = {W. H. Freeman},
  year      = {1973}
}
\end{filecontents}
\addbibresource{refs-citemsp.bib}

% ── Title ─────────────────────────────────────────────────────────────────────
\title{Draft: \textbf{\texttt{citemsp}: Citation Locators for \LaTeX}}

\author{Apostolos Tsampodimos%
  \,\,\githublink{atsampodimos/citemsp}%
  \,\,\orcidlink{0009-0001-0230-5647}}
\affil{Center for Physical Sciences and Technology\\
  Saul\.etekio~3, 10257 Vilnius, Lithuania\\
  \texttt{apostolos.tsampodimos@ftmc.lt}}

\date{\today}

\begin{document}

\maketitle

\begin{abstract}
\noindent
I present \texttt{citemsp}, a \LaTeX{} package extension of \texttt{biblatex} with per-key locators rendered as compact superscript/subscript pairs on numeric citation labels\footnote{The package is available at \url{https://github.com/apostolostsampodimos/citemsp} and on CTAN.}. The default locators are section~(\S) and paragraph~(\P), while an extensible prefix registry provides additional types -- the user can register custom prefixes with \texttt{\textbackslash citemspprefix}. The package requires no external dependencies beyond \texttt{biblatex}, \texttt{graphicx}, and \texttt{amssymb}.
\end{abstract}


\section{Syntax and Usage}

The \texttt{\textbackslash citemsp} command embeds precise locators directly in multi-source citations. It accepts section~(\S) and paragraph~(\P) locators by default, plus seven built-in prefix types -- equation~(eq.), box~($\square$), definition~($\triangleq$), figure~($\boxdot$), footnote~($\dagger$), page~(pp.), and table~($\boxplus$) -- via a slash-delimited syntax with single-letter prefixes.

The command accepts a comma-separated list of citation keys with optional locators using forward-slash delimiters. By default, the first slot renders as \S\ (section) and the second as \P\ (paragraph). A single-letter prefix changes the locator type:

\begin{tcolorbox}[colback=gray!5, colframe=gray!40, boxrule=0.5pt, arc=2pt, top=6pt, bottom=6pt, left=10pt, right=10pt]
\begin{tabular*}{\dimexpr\linewidth-0pt}{@{}l@{\extracolsep{\fill}}c@{\quad}l@{}}
\texttt{\textbackslash citemsp\{key\}}                    & $\longrightarrow$ & \citemsp{wald1984}\\[3pt]
\texttt{\textbackslash citemsp\{key/section\}}             & $\longrightarrow$ & \citemsp{wald1984/6.2}\\[3pt]
\texttt{\textbackslash citemsp\{key/section/paragraph\}}   & $\longrightarrow$ & \citemsp{wald1984/6.2/3}\\[3pt]
\texttt{\textbackslash citemsp\{key/section/e4.3\}}        & $\longrightarrow$ & \citemsp{wald1984/6.2/e4.3}\\[3pt]
\texttt{\textbackslash citemsp\{key/b21.1\}}               & $\longrightarrow$ & \citemsp{mtw1973/b21.1}\\[3pt]
\texttt{\textbackslash citemsp\{key/section/b8.1\}}        & $\longrightarrow$ & \citemsp{mtw1973/8.4/b8.1}
\end{tabular*}
\end{tcolorbox}

\noindent The seven built-in prefixes are listed in Table~\ref{tab:prefixes}, and the user can register custom prefixes with \texttt{\textbackslash citemspprefix\{letter\}\{label\}}.

\begin{table}[h]
\centering
\rowcolors{2}{blue!3}{white}
\begin{tabular}{c l l l}
\toprule
\rowcolor{blue!15}
\textbf{Prefix} & \textbf{Locator} & \textbf{Render} & \textbf{Output} \\
\midrule
\texttt{b} & Box         & $\square$     & \citemsp{mtw1973/b21.1} \\
\texttt{d} & Definition  & $\triangleq$  & \citemsp{wald1984/3.2/d5} \\
\texttt{e} & Equation    & eq.           & \citemsp{wald1984/6.2/e4.3} \\
\texttt{f} & Figure      & $\boxdot$     & \citemsp{wald1984/6.1/f3} \\
\texttt{n} & Footnote    & $\dagger$     & \citemsp{hawking1974/1/n7} \\
\texttt{p} & Page        & pp.           & \citemsp{wald1984/p42} \\
\texttt{t} & Table       & $\boxplus$    & \citemsp{mtw1973/8.4/t2} \\
\bottomrule
\end{tabular}
\caption{Registry: each single-letter prefix maps to a typographical symbol rendered before the locator number.}
\label{tab:prefixes}
\end{table}


\section{Examples}

\subsection{Single Source}

For a single reference with precise location:

\begin{tcolorbox}[colback=blue!3, colframe=blue!30, boxrule=0.5pt, arc=2pt, top=6pt, bottom=6pt, left=10pt, right=10pt]
\begin{verbatim}
The Schwarzschild metric~\citemsp{wald1984/6.1/2} describes
spherically symmetric vacuum solutions.
\end{verbatim}
\end{tcolorbox}

\noindent\textbf{Output:} The Schwarzschild metric~\citemsp{wald1984/6.1/2} describes spherically symmetric vacuum solutions.


\subsection{Multiple Sources}

For claims requiring verification across multiple independent sources:

\begin{tcolorbox}[colback=blue!3, colframe=blue!30, boxrule=0.5pt, arc=2pt, top=6pt, bottom=6pt, left=10pt, right=10pt]
\begin{verbatim}
Black hole thermodynamics~\citemsp{hawking1974/1/2, wald1984/12.5/1}
establishes a deep connection between gravity and quantum mechanics.
\end{verbatim}
\end{tcolorbox}

\noindent\textbf{Output:} Black hole thermodynamics~\citemsp{hawking1974/1/2, wald1984/12.5/1} establishes a deep connection between gravity and quantum mechanics.


\subsection{Partial Locators}

Not all entries require full locators:

\begin{tcolorbox}[colback=blue!3, colframe=blue!30, boxrule=0.5pt, arc=2pt, top=6pt, bottom=6pt, left=10pt, right=10pt]
\begin{verbatim}
General covariance~\citemsp{einstein1915, wald1984/4, hawking1974/1/2}
forms the foundation of general relativity.
\end{verbatim}
\end{tcolorbox}

\noindent\textbf{Output:} General covariance~\citemsp{einstein1915, wald1984/4, hawking1974/1/2} forms the foundation of general relativity.


\subsection{Equation Locator}

Use the \texttt{e} prefix to reference a specific equation:

\begin{tcolorbox}[colback=blue!3, colframe=blue!30, boxrule=0.5pt, arc=2pt, top=6pt, bottom=6pt, left=10pt, right=10pt]
\begin{verbatim}
The field equations~\citemsp{wald1984/3.2/e4.3} can be derived
from a variational principle.
\end{verbatim}
\end{tcolorbox}

\noindent\textbf{Output:} The field equations~\citemsp{wald1984/3.2/e4.3} can be derived from a variational principle.


\subsection{Box Locator}

Use the \texttt{b} prefix for box environments, as used extensively in Misner, Thorne \& Wheeler's \textit{Gravitation}~\cite{mtw1973}:

\begin{tcolorbox}[colback=blue!3, colframe=blue!30, boxrule=0.5pt, arc=2pt, top=6pt, bottom=6pt, left=10pt, right=10pt]
\begin{verbatim}
The stress-energy tensor~\citemsp{mtw1973/b21.1} is discussed
in a self-contained box.
\end{verbatim}
\end{tcolorbox}

\noindent\textbf{Output:} The stress-energy tensor~\citemsp{mtw1973/b21.1} is discussed in a self-contained box.


\subsection{Definition Locator}

Use the \texttt{d} prefix to reference a specific definition:

\begin{tcolorbox}[colback=blue!3, colframe=blue!30, boxrule=0.5pt, arc=2pt, top=6pt, bottom=6pt, left=10pt, right=10pt]
\begin{verbatim}
A globally hyperbolic spacetime~\citemsp{wald1984/8.1/d3}
satisfies strong causality with compact causal diamonds.
\end{verbatim}
\end{tcolorbox}

\noindent\textbf{Output:} A globally hyperbolic spacetime~\citemsp{wald1984/8.1/d3} satisfies strong causality with compact causal diamonds.


\subsection{Figure Locator}

Use the \texttt{f} prefix to point the reader to a specific figure:

\begin{tcolorbox}[colback=blue!3, colframe=blue!30, boxrule=0.5pt, arc=2pt, top=6pt, bottom=6pt, left=10pt, right=10pt]
\begin{verbatim}
The Penrose diagram~\citemsp{hawking1974/1/f2} illustrates
the causal structure near the event horizon.
\end{verbatim}
\end{tcolorbox}

\noindent\textbf{Output:} The Penrose diagram~\citemsp{hawking1974/1/f2} illustrates the causal structure near the event horizon.


\subsection{Footnote Locator}

Use the \texttt{n} prefix to cite a specific footnote in the source:

\begin{tcolorbox}[colback=blue!3, colframe=blue!30, boxrule=0.5pt, arc=2pt, top=6pt, bottom=6pt, left=10pt, right=10pt]
\begin{verbatim}
The sign convention is clarified in a
footnote~\citemsp{wald1984/3.4/n2}.
\end{verbatim}
\end{tcolorbox}

\noindent\textbf{Output:} The sign convention is clarified in a footnote~\citemsp{wald1984/3.4/n2}.


\subsection{Page Locator}

Use the \texttt{p} prefix to reference a specific page:

\begin{tcolorbox}[colback=blue!3, colframe=blue!30, boxrule=0.5pt, arc=2pt, top=6pt, bottom=6pt, left=10pt, right=10pt]
\begin{verbatim}
The ADM formalism~\citemsp{mtw1973/p520} decomposes
spacetime into spatial hypersurfaces.
\end{verbatim}
\end{tcolorbox}

\noindent\textbf{Output:} The ADM formalism~\citemsp{mtw1973/p520} decomposes spacetime into spatial hypersurfaces.


\subsection{Table Locator}

Use the \texttt{t} prefix to reference a specific table:

\begin{tcolorbox}[colback=blue!3, colframe=blue!30, boxrule=0.5pt, arc=2pt, top=6pt, bottom=6pt, left=10pt, right=10pt]
\begin{verbatim}
The Petrov classification~\citemsp{wald1984/7.3/t1} organises
the Weyl tensor into algebraic types.
\end{verbatim}
\end{tcolorbox}

\noindent\textbf{Output:} The Petrov classification~\citemsp{wald1984/7.3/t1} organises the Weyl tensor into algebraic types.


\subsection{Mixed Locator Types}

All locator types can be freely mixed in a single call:

\begin{tcolorbox}[colback=blue!3, colframe=blue!30, boxrule=0.5pt, arc=2pt, top=6pt, bottom=6pt, left=10pt, right=10pt]
\begin{verbatim}
See~\citemsp{wald1984/6.1/2, mtw1973/b21.1, hawking1974/1/e4.3}
for further details.
\end{verbatim}
\end{tcolorbox}

\noindent\textbf{Output:} See~\citemsp{wald1984/6.1/2, mtw1973/b21.1, hawking1974/1/e4.3} for further details.


\subsection{Custom Prefixes}

Users can register new prefix types with \texttt{\textbackslash citemspprefix}. For example, a cosmology paper might define a prefix for ``window function'' references:

\citemspprefix{w}{w.}

\begin{tcolorbox}[colback=blue!3, colframe=blue!30, boxrule=0.5pt, arc=2pt, top=6pt, bottom=6pt, left=10pt, right=10pt]
\begin{verbatim}
\citemspprefix{w}{w.}
The top-hat window~\citemsp{wald1984/4.2/w3} is
the simplest smoothing kernel.
\end{verbatim}
\end{tcolorbox}

\noindent\textbf{Output:} The top-hat window~\citemsp{wald1984/4.2/w3} is the simplest smoothing kernel.

\medskip
\noindent Similarly, a prefix for axioms can use a math symbol:

\citemspprefix{a}{\ensuremath{\vdash}}

\begin{tcolorbox}[colback=blue!3, colframe=blue!30, boxrule=0.5pt, arc=2pt, top=6pt, bottom=6pt, left=10pt, right=10pt]
\begin{verbatim}
\citemspprefix{a}{\ensuremath{\vdash}}
The equivalence principle~\citemsp{mtw1973/16.4/a2}
underpins all metric theories of gravity.
\end{verbatim}
\end{tcolorbox}

\noindent\textbf{Output:} The equivalence principle~\citemsp{mtw1973/16.4/a2} underpins all metric theories of gravity.


\section{Relation to Biblatex}

The standard \texttt{biblatex} provides locators via its postnote mechanism (e.g., \verb|\autocite[§3.2]{wald1984}|), but \texttt{\textbackslash citemsp} provides a compact way of guiding the reader to the precise location within the cited work, directly in the running text.


\section{Implementation}

\begin{enumerate}
  \item \textbf{Locator scaling.} \texttt{\textbackslash citemspscale} (default~0.35) and \texttt{\textbackslash citesize} control the size of locator glyphs via \texttt{\textbackslash scalebox} from \texttt{graphicx}.

  \item \textbf{Citation rendering.} A custom \texttt{biblatex} cite command, \texttt{\textbackslash citelinksp}, declared via \texttt{\textbackslash DeclareCiteCommand}, prints the hyperlinked citation number and conditionally attaches superscript/subscript locators through the dispatcher. When \texttt{hyperref} draws bordered link boxes (\texttt{colorlinks=false}), a horizontal kern prevents overlap.

  \item \textbf{Entry parsing.} The user-facing \texttt{\textbackslash citemsp} uses \texttt{\textbackslash forcsvlist} to iterate over comma-separated entries. Each entry is split at forward slashes using a delimited \texttt{\textbackslash def} pattern, extracting the citation key and up to two locator slots before passing them to the renderer.
\end{enumerate}



\printbibliography

\end{document}
